\begin{frame}{ORM -- Pros and Cons}
  \begin{columns}[T]
    \column{0.50\linewidth}
      \begin{block}{Advantages}
        \begin{itemize}
          \item \textbf{SQL Injection protection}: ORMs use parameterised
                queries internally -- injection is structurally impossible
                when used correctly.
          \item \textbf{Developer productivity}: no repetitive
                \texttt{INSERT}/\texttt{UPDATE} boilerplate; call
                \texttt{repo.save(entity)} instead.
          \item \textbf{Database portability}: switching from MySQL to
                PostgreSQL often requires only a config change.
          \item \textbf{Schema-as-code}: the DB schema is defined as Java
                classes and tracked in version control.
          \item \textbf{Automatic relationships}: \texttt{@ManyToOne}
                annotations generate correct JOINs transparently.
        \end{itemize}
      \end{block}

    \column{0.48\linewidth}
      \begin{alertblock}{Disadvantages / Trade-offs}
        \begin{itemize}
          \item \textbf{Suboptimal SQL}: generated queries can be verbose,
                inefficient, or include unnecessary JOINs.
          \item \textbf{N+1 problem}: lazy loading can trigger hundreds
                of queries unnoticed (covered in section 13e).
          \item \textbf{Learning curve}: JPA/Hibernate has many
                configuration options and subtle behaviours.
          \item \textbf{Abstraction leakage}: complex queries (window
                functions, CTEs) are hard to express in JPQL and
                require native SQL fallback.
          \item \textbf{Over-fetching}: loading full entities when only
                one column is needed wastes memory and bandwidth.
        \end{itemize}
      \end{alertblock}
  \end{columns}
\end{frame}

% ----------------------------------------------------------
\begin{frame}{ORM vs Raw SQL -- When to Use What}
  \begin{center}
  \renewcommand{\arraystretch}{1.3}
  \begin{tabularx}{\linewidth}{@{}lXX@{}}
    \toprule
    \textbf{Scenario} & \textbf{ORM} & \textbf{Raw SQL} \\
    \midrule
    Standard CRUD       & Excellent -- minimal code & Unnecessary boilerplate \\
    Simple lookups      & Great (\texttt{findBy...}) & Overkill \\
    Complex aggregations & Awkward JPQL & Better -- native SQL \\
    Bulk inserts (10k+) & Slow (entity-by-entity) & Use \texttt{LOAD DATA} or batch JDBC \\
    Window functions    & Not supported in JPQL & Required \\
    Reporting queries   & Poor -- generates many JOINs & Direct SQL or views \\
    \bottomrule
  \end{tabularx}
  \end{center}

  \begin{block}{Best Practice: Mix Both}
    Use the ORM for standard CRUD and relationship navigation.
    Drop down to native SQL (via \texttt{@Query(nativeQuery=true)} in Spring Data or
    \texttt{EntityManager.createNativeQuery()}) for performance-critical operations.
  \end{block}
\end{frame}
